\documentclass[11pt,a4paper]{article}
\usepackage[T1]{fontenc}
\usepackage{lmodern}
\usepackage{amsmath,amssymb,amsthm}
\usepackage[margin=24mm]{geometry}
\usepackage{microtype}
\usepackage[hidelinks]{hyperref}
\hypersetup{
  pdftitle={An elementary Diophantine obstruction},
  pdfsubject={A self-contained proof that the equation has no integer solutions},
  pdfkeywords={Diophantine equation, non-existence, cyclotomic divisor, elementary proof}
}
\newtheorem*{theorem}{Theorem}
\newtheorem*{lemma}{Lemma}
\setlength{\parindent}{0pt}
\setlength{\parskip}{6pt}
\setlength{\emergencystretch}{2em}
\allowdisplaybreaks[1]

\begin{document}
\begin{center}
  {\Large\bfseries An elementary obstruction for}\\[5pt]
  {\Large $y^3+y^2+xy=x^4+1$}\\[7pt]
  {\small A self-contained non-existence proof}
\end{center}
\vspace{2pt}

\begin{theorem}
There is no pair $(x,y)\in\mathbb Z^2$ satisfying
\begin{equation}\label{eq:original}
 y^3+y^2+xy=x^4+1.
\end{equation}
\end{theorem}

\textbf{An elementary divisor observation.}
\begin{lemma}\label{lem:divisors}
Let $u\in\mathbb Z$. Every positive divisor $d$ of $u^2+u+1$ with
$3\nmid d$ satisfies $d\equiv1\pmod3$.
\end{lemma}
\begin{proof}
Let $p\ne3$ be a prime divisor of $u^2+u+1$. Modulo $p$, we have
$u^3=1$ and $u\ne1$: the latter follows because substitution of $u=1$
would give $p\mid3$. Multiplication by $u$ therefore partitions the
$p-1$ nonzero residue classes into disjoint triples
$\{r,ur,u^2r\}$. These three classes are distinct, since $u^3=1$ and
$u\ne1$. Hence $3\mid p-1$. Every prime factor of $d$ is consequently
$1\pmod3$, and so is $d$; this includes the empty product $d=1$.
\end{proof}

\textbf{Proof of the theorem.}
Suppose \eqref{eq:original} holds. Reduction modulo $3$ gives
\[
 0\equiv y^2+(x+1)y-x^2-1
   \equiv (x-1)^2+(x-y+1)^2\pmod3.
\]
The only squares modulo $3$ are $0$ and $1$, so both squares vanish. Thus
\begin{equation}\label{eq:residues}
 x\equiv1,\qquad y\equiv2\pmod3.
\end{equation}
Define
\[
 A=x^2+y,\qquad
 B=x^4-x^2y+y^2-x^2+y+x.
\]
A direct expansion, followed by \eqref{eq:original}, yields
\begin{equation}\label{eq:factorization}
 AB=x^6+x^3+1
   =\left(x^3+\frac12\right)^2+\frac34>0.
\end{equation}
In particular $A\ne0$. If $A<0$, then $t=-y\ge x^2+1$ and
\[
 y^2+y+x=t(t-1)+x
   \ge x^2(x^2+1)+x=x^4+x(x+1)\ge0,
\]
where $x(x+1)\ge0$ for every integer $x$. Since $y<0$, this contradicts
$y(y^2+y+x)=x^4+1>0$. Therefore $A>0$, and \eqref{eq:factorization}
also gives $B>0$.

On the other hand, \eqref{eq:residues} implies
\[
 B\equiv1-2+1-1+2+1\equiv2\pmod3.
\]
But $B$ is a positive divisor of $(x^3)^2+x^3+1$ not divisible by $3$.
The lemma, with $u=x^3$, forces $B\equiv1\pmod3$, a contradiction.
\hfill$\square$

\newpage
\section*{How the obstruction is found}

\textbf{Make two terms cancel before choosing a modulus.}
Write
\[
 F(x,y)=y^3+y^2+xy-x^4-1.
\]
The terms $y^2$ and $x^4$ suggest substituting $y=-x^2$. This cancels
those two terms exactly and leaves
\[
 F(x,-x^2)=-(x^6+x^3+1).
\]
Consequently, an integer zero of $F$ must make $x^2+y$ a divisor of
$x^6+x^3+1$. The useful modulus is therefore not a fixed small integer:
it is the expression $x^2+y$ suggested by the equation itself.

\textbf{Recognize the restriction on prime divisors.}
The remaining expression is
\[
 x^6+x^3+1=u^2+u+1,\qquad u=x^3.
\]
A prime divisor other than $3$ must admit a nontrivial cube root of
unity, which is exactly what the triples in the lemma express.
It must therefore be $1\pmod3$. This turns the substitution into a
restriction on \emph{every positive divisor}, rather than merely a
congruence on $x$ or $y$.

\textbf{Use the complementary factor.}
The original equation modulo $3$ forces $(x,y)\equiv(1,-1)$, so
$x^2+y\equiv0\pmod3$. The exceptional prime $3$ thus obscures the
obstruction in this first factor. Polynomial division isolates the
complementary factor $B$ through the exact identity
\[
 F(x,y)=(x^2+y)
 \bigl(x^4-x^2y+y^2-x^2+y+x\bigr)-(x^6+x^3+1).
\]
Unlike $x^2+y$, this complementary factor is $-1\pmod3$, so it is
coprime to $3$ and directly contradicts the lemma once it is positive.

\textbf{Do not omit the sign check.}
A negative divisor can be $2\pmod3$ without having any prime factor
$2\pmod3$; for example, $-1$ does this. Thus the positivity argument is
essential. The equation itself rules out $x^2+y<0$, and the positive
product in \eqref{eq:factorization} then forces $B>0$. No sign assumption
on the unknown integers is being made.

The proof combines the rigid residue pair modulo $3$ with a
prime-divisor restriction exposed by cancellation. It excludes every
integer pair, without a search bound or an unproved auxiliary claim.

\end{document}
